% Compile: xelatex cover_letter.tex
% Self-contained: no images or other local assets are required.
% RSOI is described as ongoing work; no completed Go deployment is asserted.
\documentclass[11pt,a4paper]{article}

\usepackage{fontspec}
\usepackage{polyglossia}
\setmainlanguage{russian}
\setotherlanguage{english}
\setmainfont{Liberation Sans}
\usepackage[left=2.2cm,right=2.2cm,top=1.8cm,bottom=1.8cm]{geometry}
\usepackage{xcolor}
\usepackage{hyperref}
\usepackage{microtype}
\usepackage{setspace}

\definecolor{accent}{HTML}{2B4C7E}
\definecolor{ink}{HTML}{28313D}
\definecolor{muted}{HTML}{657080}
\definecolor{rulecolor}{HTML}{CAD4E1}

\hypersetup{
  colorlinks=true,
  urlcolor=accent,
  linkcolor=accent,
  pdftitle={Сопроводительное письмо --- Алексеев Даниил},
  pdfauthor={Алексеев Даниил Александрович},
  pdfsubject={Стажировка в backend-разработке}
}
\pagestyle{empty}
\setlength{\parindent}{0pt}
\setlength{\parskip}{0.75em}
\setlength{\emergencystretch}{2em}
\setstretch{1.08}
\raggedright

\begin{document}
\color{ink}

% ---------- Personal header ----------
{\small\color{accent}КАНДИДАТ НА СТАЖИРОВКУ \enspace / \enspace BACKEND}\par
\vspace{0.1em}
{\fontsize{23}{27}\selectfont\bfseries Алексеев Даниил}\par
\vspace{-0.1em}
{\small\color{muted}
  Москва \enspace\textbullet\enspace
  \href{mailto:dalyoksi@yandex.ru}{dalyoksi@yandex.ru}
  \enspace\textbullet\enspace +7\,911\,420-29-96
}\par
\vspace{0.3em}
{\color{accent}\rule{\textwidth}{1.2pt}}\par
\vspace{0.5em}

{\fontsize{16}{20}\selectfont\bfseries Сопроводительное письмо}\par
\vspace{0.25em}

Здравствуйте!

Меня зовут Даниил Алексеев. Учусь в магистратуре МГТУ им. Н.Э. Баумана по направлению «Программная инженерия» и хочу развиваться в backend-разработке.

\textbf{Больше всего меня интересуют компьютерные сети и распределённые системы.} Мне нравится разбираться, как запрос проходит между компонентами приложения, что происходит при недоступности одного из сервисов и как организовать восстановление работы. Особенно интересны межсервисное взаимодействие, асинхронная обработка запросов и согласованность данных.

Хочу заниматься разработкой сервисов и инфраструктурных компонентов: проектировать API, работать с хранилищами и очередями, исследовать причины ошибок и улучшать производительность. Мне важно понимать связь между кодом приложения, сетевым взаимодействием и поведением операционной системы.

Моя основа - C++ и олимпиадное программирование: дважды становился призёром полуфинала ICPC. В командном грантовом проекте разрабатывал REST API, работал с PostgreSQL и контейнеризацией сервисов. Во время оплачиваемой учебной практики создал приложение на Python для расчётов по STL/DXF-моделям с графическим интерфейсом и 3D-визуализацией.

Сейчас работаю над микросервисным проектом на Go в рамках курса «Распределённые системы обработки информации». В нём разбираюсь с обработкой частичных отказов, повторным выполнением запросов и развёртыванием сервисов. Пятый год использую Linux как desktop ОС - Debian и дистрибутивы на базе Ubuntu; углубляю знания сетей через видеолекции и практические задания в eNSP.

На стажировке хочу применить эти знания к задачам реального продукта и пройти полный цикл разработки: от обсуждения решения до тестирования и выпуска. Для меня важны обратная связь по коду, возможность задавать вопросы опытным коллегам и постепенно брать на себя больше ответственности за результат.

Буду рад обсудить задачи команды и подробнее рассказать о своих проектах.

\vspace{0.3em}
С уважением,\\[0.25em]
\textbf{Даниил Алексеев}


\end{document}
